% ---
% id: reductions/rem:two-sided-scope
% kind: rem
% title: "Scope of the two-sided reduction"
% summary: >-
%   The two-sided reduction relates no marginals of the specified and virtual
%   descriptions and licenses no entry whose right-hand side is not a real
%   constant. Its order-independence concerns soundness; the order a chain
%   records stays data.
% prominence: standard
% status: accepted
% uses: [thm:two-sided]
% mentions: [rem:source-transfer-scope, lem:precomposition, def:chain, rem:ordering]
% assumes: []
% refs: []
% provenance: newly-derived
% lean:
% ---

\begin{remark}\label{rem:two-sided-scope}
Theorem~\ref{thm:two-sided} uses and implies no relation between a marginal
of the specified description and one of the virtual description, and none
would follow (Remark~\ref{rem:source-transfer-scope}). It licenses no entry
whose right-hand side is not a real constant:
Lemma~\ref{lem:precomposition}\,(iv) is a counterexample for such an entry,
and the admission rule of Definition~\ref{def:chain} refuses one beside a
non-empty source segment. And part (i) does not say that a chain may be
\emph{recorded} in either order: the record's order is data, stated as
condition (C3) of Definition~\ref{def:chain}, and what (i) says is that for a
licensed list the order is not what carries the soundness
(Remark~\ref{rem:ordering}).
\end{remark}
